\subsection{Case I: a proper rigid subgraph}
\label{sec:molecular-algebraic-induction-caseI}

This subsection proves Katoh--Tanigawa's Case~I: a minimal $k$-dof-graph
with a proper rigid subgraph $H$ realizes its target rank by contracting
$H$ to a single pinned body and gluing the rigid block to the
contraction's own inductive realization (KT eq.~(6.3)). Two pieces of
shared infrastructure come first: the panel-cycle realization (KT
Lemma~5.4), used later in Case~III, and the block pin --- the
framework-side model of contracting a rigid block to one body. The Case~I
accounting itself (\cref{lem:case-I}) is genericity-free; genericity
enters only through KT Claim~6.4 (\cref{lem:claim-6-4}), needed to realize
both legs of the block-triangular splice at a single common placement.
The realization lemmas (\cref{lem:case-I-realization}, \cref{lem:case-I-realization-nonsimple},
\cref{lem:case-I-realization-h65}, \cref{lem:case-I-dispatch}) assemble
this into KT Lemmas~6.2, 6.3, and 6.5.

Katoh--Tanigawa's Lemma~5.4 realizes a cycle graph directly: for
$3 \le m \le D$, a cycle of $m$ bodies whose consecutive panels have
linearly independent supporting extensors is infinitesimally rigid at the
full rank $D(m-1)$. The argument has four pieces: the two-body base case;
the dimension bound $m \le D$ forced by independence in the
$D$-dimensional screw space; the telescoping argument generalizing the
base case to any cycle length; and the existence of $m$ independent
panels for $3 \le m \le D$, an exterior-algebra statement on the
grade-$2$ joins of the panel normals that a basis choice on
$\Lambda^2 \R^{k+2}$ settles.

\begin{lemma}[Cycle realization, short-cycle base]
  \label{lem:cycle-realization-base}
  \lean{CombinatorialRigidity.Molecular.PanelHingeFramework.toBodyHinge_rankHypothesis_zero}
  \leanok
  \uses{def:panel-hinge-framework, def:rank-hypothesis, lem:theorem-55-base}
  A panel-hinge framework $P$ with two edges $e_1, e_2$ joining two
  distinct bodies $u \ne v$, whose panel support extensors
  $C(p(e_1)), C(p(e_2))$ are linearly independent, has a body-hinge
  interpretation infinitesimally rigid on $\{u,v\} = V(G)$
  (\cref{def:rank-hypothesis}), the full rank $D(2-1) = D$.
\end{lemma}
\begin{proof}
  \leanok
  \uses{lem:theorem-55-base}
  The body-hinge interpretation carries exactly the two independent
  supporting extensors as its hinge data, so the parallel-hinges-full base
  case (\cref{lem:theorem-55-base}) applies verbatim.
\end{proof}

\begin{lemma}[Rigid cycle dimension bound]
  \label{lem:cycle-realization-dim-bound}
  \lean{CombinatorialRigidity.Molecular.card_le_screwDim_of_linearIndependent,
        CombinatorialRigidity.Molecular.PanelHingeFramework.card_le_screwDim_of_supportExtensor_linearIndependent}
  \leanok
  \uses{def:panel-hinge-framework, def:rigidity-matrix}
  If the supporting extensors of $m$ edges of a panel-hinge framework are
  linearly independent in the screw space $\Lambda^k \R^{k+2}$, then
  $m \le D = \binom{k+2}{2}$.
\end{lemma}
\begin{proof}
  \leanok
  A linearly independent family in a $D$-dimensional space has at most $D$
  members, and $\dim \Lambda^k \R^{k+2} = D$ (\cref{def:rigidity-matrix}).
\end{proof}

\begin{lemma}[Panel independence reduces to grade-$2$ join independence]
  \label{lem:panel-support-extensor-independence}
  \lean{CombinatorialRigidity.Molecular.panelSupportExtensor_linearIndependent_iff}
  \leanok
  \uses{def:panel-support-extensor, def:meet-complement-iso}
  A family of $m$ panel support extensors
  $i \mapsto C(p(e_i)) = C(n_{u_i}, n_{v_i})$ is linearly independent in
  the screw space $\Lambda^k \R^{k+2}$ iff the family of grade-$2$ joins
  $i \mapsto n_{u_i} \wedge n_{v_i}$ is independent in
  $\Lambda^2 \R^{k+2}$.
\end{lemma}
\begin{proof}
  \leanok
  The complement iso $\Lambda^2 \R^{k+2} \simeq \Lambda^k \R^{k+2}$
  (\cref{def:meet-complement-iso}) is a linear equivalence, and the panel
  support extensor is its image of the grade-$2$ join by definition
  (\cref{def:panel-support-extensor}); a linear equivalence both
  preserves and reflects linear independence.
\end{proof}

\begin{lemma}[Existence of independent panels for $m \le D$]
  \label{lem:exists-independent-panel-extensor}
  \lean{CombinatorialRigidity.Molecular.exists_independent_panelSupportExtensor,
        CombinatorialRigidity.Molecular.exists_independent_normalsJoin}
  \leanok
  \uses{def:panel-support-extensor, lem:panel-support-extensor-independence,
        lem:extensor-independence}
  For every $m \le D = \binom{k+2}{2}$ there are $m$ pairs of panel
  normals $(n_{u_i}, n_{v_i})$ whose supporting extensors are linearly
  independent in the screw space $\Lambda^k \R^{k+2}$.
\end{lemma}
\begin{proof}
  \leanok
  By \cref{lem:panel-support-extensor-independence} it suffices to
  produce $m$ independent grade-$2$ joins. The index set of $2$-element
  subsets of $\{1, \dots, k+2\}$ has cardinality $\binom{k+2}{2} = D \ge
  m$; picking $m$ distinct $2$-subsets and taking the corresponding
  pairs of standard basis vectors gives $m$ joins forming a subfamily of
  the basis-indexed exterior-power family of $\Lambda^2 \R^{k+2}$, which
  is linearly independent by the extensor-independence Lemma~2.1
  (\cref{lem:extensor-independence}); a subfamily of an independent
  family is independent.
\end{proof}

\begin{lemma}[Rigidity of a cycle with independent extensors]
  \label{lem:cycle-realization-rigid}
  \lean{CombinatorialRigidity.Molecular.PanelHingeFramework.toBodyHinge_rankHypothesis_zero_cycle,
        CombinatorialRigidity.Molecular.BodyHingeFramework.rankHypothesis_zero_of_cycle,
        CombinatorialRigidity.Molecular.BodyHingeFramework.isTrivialMotion_of_isInfinitesimalMotion_cycle,
        CombinatorialRigidity.Molecular.BodyHingeFramework.theorem_55_cycle}
  \leanok
  \uses{def:panel-hinge-framework, def:rank-hypothesis, lem:cycle-realization-base}
  A panel-hinge framework on the cycle of $m$ bodies $0, 1, \dots, m-1$
  --- the $i$-th hinge joining body $i$ to body $i+1$ (cyclically) ---
  whose $m$ supporting extensors are linearly independent in the screw
  space $\Lambda^k \R^{k+2}$ is infinitesimally rigid, realizing the rank
  hypothesis at $k' = 0$ (\cref{def:rank-hypothesis}). The same conclusion
  holds when the cycle is indexed by an arbitrary map
  $\mathrm{vtx} \colon \Z/m \to V$, with no injectivity required: rigidity
  is then asserted on the bodies the cycle actually visits,
  $\operatorname{range}(\mathrm{vtx})$.
\end{lemma}
\begin{proof}
  \leanok
  \uses{lem:rank-parallel-full}
  Each hinge constraint puts the relative screw $S(i) - S(i+1)$ of an
  infinitesimal motion $S$ in the one-dimensional span of its supporting
  extensor, and the $m$ differences telescope around the cycle to
  $\sum_i (S(i) - S(i+1)) = 0$. Independence of the $m$ extensors then
  forces every difference to vanish --- the $m$-edge generalization of
  the parallel-hinges-full Lemma~5.3 (\cref{lem:rank-parallel-full}, the
  $m=2$ case) --- so $S$ is constant around the connected cycle, hence a
  trivial motion, and \cref{def:rank-hypothesis} gives the rank hypothesis
  at $k'=0$.
\end{proof}

\begin{lemma}[Cycle realization; KT Lemma 5.4]
  \label{lem:cycle-realization}
  \lean{CombinatorialRigidity.Molecular.PanelHingeFramework.cycle_realization}
  \leanok
  \uses{def:rigidity-matrix, def:rank-hypothesis, def:cycle-data, def:k-dof,
        lem:cycle-realization-base, lem:cycle-realization-dim-bound,
        lem:cycle-realization-rigid, lem:cycle-normals,
        lem:case-III-claim612-line-in-panel-union, lem:genericity-device}
  A minimal $0$-dof-graph $G$ that \emph{is} a cycle graph
  (\cref{def:cycle-data}) on $m$ vertices, $3 \le m \le D$, has a generic
  full-rank panel-hinge realization $(G,p)$ (\cref{def:rank-hypothesis}):
  general-position, algebraically independent panel normals attaining the
  full rank $D(|V|-1)$.
\end{lemma}
\begin{proof}
  \leanok
  \uses{lem:cycle-realization-rigid,
        lem:cycle-normals, lem:case-III-claim612-line-in-panel-union}
  The cyclic shared-normal seed (\cref{lem:cycle-normals}) supplies $m$
  panel normals --- one per body, consecutive hinges sharing a panel ---
  whose $m$ cyclic supporting extensors are linearly independent and
  whose cyclic grade-$2$ joins are nonzero ($m \le d = k+1$, the
  dimension chain from $D = \binom{d+1}{2} = \binom{k+2}{2}$). Assign
  these normals to the cycle vertices and orient each edge by the
  graph's canonical endpoint choice: every cycle edge's supporting
  extensor is then $\pm$ a member of the seed's independent cyclic
  family, so the edge-extensor family is independent and nonvanishing
  edge-by-edge (every edge of $G$ \emph{is} a cycle edge). The
  telescoping rigidity (\cref{lem:cycle-realization-rigid}, in its
  arbitrary-body-type form) makes this realization infinitesimally rigid
  on the cycle bodies, i.e.\ on $V(G)$, and the bare-to-generic upgrade
  (\cref{lem:case-III-claim612-line-in-panel-union}) re-realizes it
  generically at full rank, exactly as for the triangle base
  (\cref{lem:triangle-realization}).
\end{proof}

\begin{fmlnote}
  KT states Lemma~5.4 for the full range $3 \le m \le D$; the formalized
  statement covers $m \le d$ (with $D = \binom{d+1}{2}$), the range the
  chain-extraction dichotomy of \cref{sec:molecular-algebraic-induction-caseIII}
  actually consumes --- the remaining band $d < m \le D$ has no consumer in
  this induction. Only the ``independent $\Rightarrow$ rigid'' direction of
  the Crapo--Whiteley equivalence is used (KT realizes the cycle \emph{at} an
  independent hinge family); the projective assembly of the realization on a
  concrete cycle graph, which KT cites rather than reproves, is also carried
  out here. This is Katoh--Tanigawa's Lemma~5.4; its geometric content is due
  to Crapo--Whiteley~\cite[Proposition~3.4]{crapoWhiteley1982} (the $d=3$
  case) and Whiteley~\cite[Proposition~3]{whiteley1999}, the argument
  extending to general $d$ without modification.
\end{fmlnote}

The block pin is the framework-side model of contracting a set of bodies
to a single pinned one, the move Case~I makes on a rigid subgraph.

\begin{definition}[Block-pinning a set of bodies]
  \label{def:pinned-motions-on}
  \lean{CombinatorialRigidity.Molecular.BodyHingeFramework.pinnedMotionsOn,
        CombinatorialRigidity.Molecular.BodyHingeFramework.pinnedMotionsOn_singleton,
        CombinatorialRigidity.Molecular.BodyHingeFramework.pinnedMotionsOn_eq_iInf}
  \leanok
  \uses{def:rigidity-matrix, lem:rank-delete-vertex}
  For a body-hinge framework $F = (G,p)$ and a set of bodies
  $s \subseteq V$, the \emph{block pin} $Z_s(G,p)$ is the subspace of
  infinitesimal motions vanishing on every body of $s$:
  $\{S \in Z(G,p) : S|_s = 0\}$.
\end{definition}

At a singleton $s = \{v\}$ the block pin recovers the pin-a-body subspace
$\operatorname{pinned}(v)$ of \cref{lem:rank-delete-vertex}, and for a
nonempty block it equals the intersection of the single-body pins,
$Z_s(G,p) = \bigsqcap_{v \in s} \operatorname{pinned}(v)$. Case~I applies
the block pin to a rigid subgraph $H$: pinning all of $V(H)$ to the zero
screw is the algebraic effect of contracting $H$ to a single body.

\begin{lemma}[Edge deletion enlarges the block pin]
  \label{lem:pinned-motions-on-mono}
  \lean{CombinatorialRigidity.Molecular.BodyHingeFramework.pinnedMotionsOn_le_withGraph_of_le,
        CombinatorialRigidity.Molecular.BodyHingeFramework.finrank_pinnedMotionsOn_le_of_graph_le}
  \leanok
  \uses{def:pinned-motions-on, def:framework-with-graph, lem:motions-mono-of-graph-le}
  Replacing $G$ by any subgraph $G' \le G$ (keeping the hinge data,
  \cref{def:framework-with-graph}) can only grow the block pin:
  $Z_s(G,p) \subseteq Z_s(G',p)$, hence
  $\dim Z_s(G,p) \le \dim Z_s(G',p)$.
\end{lemma}
\begin{proof}
  \leanok
  \uses{lem:motions-mono-of-graph-le}
  The block pin $Z_s(G,p)$ is $Z(G,p)$ intersected with the
  graph-independent vanishing conditions $S|_s = 0$, so the inclusion is
  the motion-space monotonicity (\cref{lem:motions-mono-of-graph-le})
  applied to the first factor; the dimension bound follows.
\end{proof}

This is the direction Case~I's block-triangular gluing travels: placing
the contraction realization on the smaller graph $G/E(H)$ and re-adding
the edges $E(H)$ only grows the block-pinned rank, the slack in
\cref{lem:pinned-motions-on-rank-bound} filled by the contraction's own
rank.

\begin{lemma}[Block-pinning drops at least the $D$ trivial-motion dimensions]
  \label{lem:pinned-motions-on-rank-bound}
  \lean{CombinatorialRigidity.Molecular.BodyHingeFramework.screwDim_add_finrank_pinnedMotionsOn_le}
  \leanok
  \uses{def:pinned-motions-on, lem:rank-delete-vertex, lem:trivial-motions-rank-bound}
  For a nonempty block of bodies $s \subseteq V$,
  \[
    D + \dim Z_s(G,p) \;\le\; \dim Z(G,p).
  \]
\end{lemma}
\begin{proof}
  \leanok
  \uses{lem:rank-delete-vertex}
  The block pin $Z_s(G,p)$ sits inside the single-body pin
  $\operatorname{pinned}(v)$ of any $v \in s$, which meets the
  $D$-dimensional trivial motions only at $0$
  (\cref{lem:rank-delete-vertex}); so the block pin and the trivial
  motions are disjoint submodules of $Z(G,p)$, and their dimensions add
  to at most $\dim Z(G,p)$.
\end{proof}

This is the lower-bound half of Case~I's block-triangular accounting: a
single-body pin is an exact $+D$ direct-sum split (\cref{lem:rank-delete-vertex}),
whereas pinning a whole rigid block $H$ collapses it to one effective
body and the residual $D(|V(H)|-1)$ constraints make the bound an
inequality --- the contraction's own rank, supplied by the induction
hypothesis, fills the slack to equality. The next lemma records the two
extreme cases of that slack.

\begin{lemma}[The block pin vanishes exactly where the framework is rigid]
  \label{lem:pinned-motions-on-rigid-iff}
  \lean{CombinatorialRigidity.Molecular.BodyHingeFramework.pinnedMotionsOn_eq_bot_of_isInfinitesimallyRigid,
        CombinatorialRigidity.Molecular.BodyHingeFramework.finrank_pinnedMotionsOn_eq_zero_of_isInfinitesimallyRigid,
        CombinatorialRigidity.Molecular.BodyHingeFramework.isInfinitesimallyRigid_of_block_of_pinnedMotionsOn_eq_bot}
  \leanok
  \uses{def:pinned-motions-on}
  If $F$ is infinitesimally rigid, then block-pinning any nonempty body
  set $s$ leaves nothing: $Z_s(G,p) = 0$. Conversely, if every
  infinitesimal motion of $F$ is constant on a nonempty block $s$ and the
  block pin vanishes, $Z_s(G,p) = 0$, then $F$ is infinitesimally rigid.
\end{lemma}
\begin{proof}
  \leanok
  A block-pinned motion is an infinitesimal motion, so rigidity makes it
  trivial (constant); vanishing at one body of $s$ then forces the
  constant to be $0$. Conversely, given a motion $S$ and a body
  $v \in s$, the constant assignment equal to $S(v)$ everywhere is a
  trivial motion; their difference vanishes on all of $s$ by the
  block-constancy hypothesis, hence lies in the (vanishing) block pin,
  forcing $S$ to equal that constant.
\end{proof}

\begin{lemma}[Case I: rigid-subgraph contraction realizes the rank; KT Lemmas 6.2, 6.3, 6.5]
  \label{lem:case-I}
  \lean{CombinatorialRigidity.Molecular.BodyHingeFramework.isInfinitesimallyRigidOn_iff_pinnedMotionsOn_le,
        CombinatorialRigidity.Molecular.PanelHingeFramework.toBodyHinge_rankHypothesis_iff_finrank_pinnedMotionsOn,
        CombinatorialRigidity.Molecular.BodyHingeFramework.rankHypothesis_iff_finrank_pinnedMotionsOn}
  \leanok
  \uses{thm:minimal-kdof-reduction, lem:contraction-minimality,
        lem:rank-delete-vertex, def:pinned-motions-on,
        lem:pinned-motions-on-rank-bound, lem:genericity-device}
  Let $G$ be a minimal $k$-dof-graph with a proper rigid subgraph $H$ on
  the (nonempty) body set $s = V(H)$. Contracting $H$ to a single pinned
  body is the block pin $Z_s(G,p)$ (\cref{def:pinned-motions-on}), the
  framework-side carrier of the contraction $G/E(H)$ (a smaller minimal
  $k$-dof-graph by \cref{lem:contraction-minimality}, supplied the
  induction hypothesis). Then $(G,p)$ realizes the $V(G)$-relative target
  rank at $k'$ (\cref{def:rank-hypothesis}) \emph{iff} the rigid block $H$
  contributes its full $D(|V(H)|-1)$ rows and the residual contraction
  contributes the rest --- the block-triangular rank addition of
  KT~(6.3): $\operatorname{rank} R(G,p) = D(|V(H)|-1) +
  \operatorname{rank} R(G/E(H), p_2)$, glued via the pin-a-body
  Lemma~5.1 (\cref{lem:rank-delete-vertex}).
\end{lemma}
\begin{proof}
  \leanok
  \uses{lem:rank-delete-vertex, lem:pinned-motions-on-rank-bound,
        lem:genericity-device}
  Pin the rigid block $H$ and place the contraction $G/E(H)$ at its
  inductive rank. The block pin and the $D$-dimensional trivial motions
  are disjoint, giving the genericity-free lower bound
  (\cref{lem:pinned-motions-on-rank-bound}). The combined rigidity matrix
  is block-triangular --- the $H$-rows occupy only the $V(H)$ columns,
  the contraction rows the rest --- so at a generic point the ranks add
  and the reverse bound holds (the entries are polynomials in
  algebraically independent panel coordinates, so a generic point attains
  the maximal rank, KT's Claim~6.4, \cref{lem:genericity-device}). The
  two bounds together pin the parent rank to $D(|V(G)|-1) - k$. This
  accounting depends only on $s, t \subseteq V(G)$, not on the ambient
  body type, so it composes through the vertex-reducing induction; the
  companion nullity form (relating $\dim Z(G,p)$ to the ambient body
  type) remains available for the deficiency and Proposition~1.1 path.
\end{proof}

\begin{lemma}[Panel-coordinate row coordinatization; KT Claim 6.4]
  \label{lem:rows-polynomial-in-normals}
  \lean{CombinatorialRigidity.Molecular.PanelHingeFramework.exists_good_realization_ofParam}
  \leanok
  \uses{def:panel-support-extensor, def:rigidity-matrix, lem:genericity-device}
  Regard a panel-hinge realization of a fixed graph $G$ as a point of
  panel-coordinate space: an assignment $q \colon V_\alpha \times \{1,
  \dots, k+2\} \to \R$ of a normal $n_v = q(v, -) \in \R^{k+2}$ to each
  body. Then the rigidity rows of $R(G,q)$ form a finite family of
  functionals on $\R^{D|V_\alpha|}$ whose coordinates, in the standard
  basis, are degree-$2$ polynomials in the entries of $q$ --- exactly the
  input shape the genericity device (\cref{lem:genericity-device}) needs
  to be applied to a \emph{varying} panel realization, rather than a
  constant one.
\end{lemma}
\begin{proof}
  \leanok
  A coordinate of the grade-$2$ join $n_u \wedge n_v$ against the
  standard dual basis of $\Lambda^2 \R^{k+2}$ is the $2 \times 2$ minor
  $n_u(i)\,n_v(j) - n_u(j)\,n_v(i)$, a determinant of the two normals'
  coordinates and hence bilinear, so it is a degree-$2$ polynomial in the
  panel coordinates $q$. The supporting extensor $C(p(e)) =
  \mathrm{complementIso}(n_u \wedge n_v)$ (\cref{def:panel-support-extensor})
  is a fixed linear image of the join, so its coordinates stay degree-$2$.
  Each rigidity row spans the annihilator of its edge's supporting
  extensor in $(\R^{k+2})^*$, and its coordinates are in turn a fixed
  linear combination of the extensor's coordinates at the two endpoints,
  so they too are degree-$2$ in $q$. Applying the genericity device to
  this family, with the witnessed independent subfamily supplied by the
  general-position seed (\cref{lem:exists-independent-panel-extensor}),
  produces a normal assignment attaining the codimension bound
  $\#s + \dim Z(G,q) \le D|V_\alpha|$ at the same corank.
\end{proof}

\begin{lemma}[Case I splice glue; KT eq.\ (6.3)]
  \label{lem:case-I-splice-seed}
  \lean{CombinatorialRigidity.Molecular.BodyHingeFramework.isInfinitesimallyRigidOn_of_splice}
  \leanok
  \uses{def:framework-with-graph, def:rank-hypothesis, lem:case-I}
  Let $H$ (on $s_H = V(H)$) be a proper rigid subgraph and $G/E(H)$ (on
  $s_c$) its contraction, both subgraphs of $G$ sharing the contracted
  body $c \in s_H \cap s_c$ and covering $V(G) \subseteq s_H \cup s_c$.
  If $F$ restricted to $H$ (\cref{def:framework-with-graph}) is
  infinitesimally rigid on $s_H$, and $F$ restricted to $G/E(H)$ is
  infinitesimally rigid on $s_c$, then $F$ is infinitesimally rigid on
  $V(G)$.
\end{lemma}
\begin{proof}
  \leanok
  \uses{lem:motions-mono-of-graph-le}
  Each leg transports to the parent by the monotonicity that re-adding
  edges only shrinks the motion space (\cref{lem:motions-mono-of-graph-le},
  since $H, G/E(H) \le G$), so a parent motion is already constant on
  $s_H$ and on $s_c$. The two constants agree at the shared contracted
  body $c$, so every motion is constant across $s_H \cup s_c \supseteq
  V(G)$.
\end{proof}

This is KT's block-triangular splice (eq.~(6.3)) read as two overlapping
rigid pieces gluing along a shared body, and it is genericity-free: the
genericity device (\cref{lem:genericity-device}) enters only at the
realization lemma \cref{lem:case-I-realization}, where the two legs must be
realized at a single generic placement --- the intersection of the two
legs' Zariski-open rigid loci. The hypotheses here are the satisfiable
inductive facts, relative rigidity of each piece on a common placement
$F$, not the parent rank they conclude.

\begin{lemma}[Graph-level minimality of the rigid-subgraph contraction]
  \label{lem:rigidContract-isMinimalKDof}
  \lean{Graph.rigidContract_isMinimalKDof}
  \leanok
  \uses{lem:contraction-minimality, def:rigid-contraction, lem:reduction-step}
  Contracting a proper rigid subgraph $H$ of a minimal $k$-dof-graph $G$
  again yields a minimal $k$-dof-graph: $(G/E(H))$ is minimal $k$-dof,
  with strictly fewer vertices (\cref{lem:reduction-step}).
\end{lemma}
\begin{proof}
  \leanok
  \uses{lem:contraction-minimality, thm:def-eq-corank, lem:two-edge-conn}
  The reduction theorem (\cref{thm:minimal-kdof-reduction}) supplies only
  the existence of a rigid subgraph together with the induction
  hypothesis for the reduction, leaving the contraction's own minimality
  to be established separately; that is the content here. The matroid
  contraction $M(\tilde G)/E(\tilde H)$ is already a minimal $k$-dof
  matroid (\cref{lem:contraction-minimality}), and the matroid of the
  \emph{graph} contraction $G/E(H)$ equals that matroid contraction ---
  the graph deletes the edges $E(H)$ and relabels vertices, while the
  matroid contracts the fibers $E(\tilde H)$; reconciling the two needs
  $\tilde H$ connected on $V(H)$, which follows from $H$ being $0$-dof
  (\cref{lem:two-edge-conn}). With that identification the base/fiber
  condition transports directly (an edge of $G/E(H)$ is a surviving
  $G$-edge), and the deficiency transports via the def-equals-corank
  bridge (\cref{thm:def-eq-corank}) against the conserved matroid rank and
  the collapse vertex-count $|V(G/E(H))| = |V(G)|-|V(H)|+1$.
\end{proof}

\begin{lemma}[KT Claim 6.4: the surviving block attains its rank at a generic placement (eq.\ (6.9))]
  \label{lem:claim-6-4}
  \lean{CombinatorialRigidity.Molecular.PanelHingeFramework.rigidContract_exterior_rank_transport_htransport}
  \leanok
  \uses{def:rigidity-matrix, def:rigid-contraction, lem:rank-delete-vertex,
        lem:contraction-minimality}
  Let $G$ be a simple minimal $0$-dof-graph with a proper rigid subgraph
  $H$ on $s = V(H)$, and write $s_c = (V(G) \setminus V(H)) \cup \{r\}$
  for the surviving body set with representative $r$. Given the
  contraction $G/E(H)$ realized at its inductive rank, the surviving-edge
  block of the rigidity matrix --- restricted to the edges $E(G)
  \setminus E(H)$ and to the columns outside $V(H)$ --- attains rank
  $D(|s_c|-1)$ at a generic placement: at a Zariski-open locus of panel
  coordinates, the surviving rows are linearly independent after
  projecting away the columns of $V(H)$ (KT's bottom-right block,
  eq.~(6.9); the row-count companion of the pin-a-body Lemma~5.1,
  \cref{lem:rank-delete-vertex}).
\end{lemma}
\begin{proof}
  \leanok
  \uses{lem:rank-delete-vertex, lem:contraction-minimality}
  The argument transports the contraction's rank across the collapse map
  $f$ sending every body of $V(H)$ to $r$ and fixing the rest, from a
  generic realization of the abstract relabelled graph $G/E(H)$ to the
  surviving-edge block of the parent.

  First, the contraction's generic realization is transported to the
  endpoint selector recording the links of $G/E(H)$ through $f$: since
  both selectors record the same links, they agree up to swap and the two
  realizations have the same rigidity and motion space. This gives a
  free-normal panel framework on $G/E(H)$ in general position, rigid on
  its whole vertex set $s_c$.

  Second --- the genuine content of the claim --- the exterior-column
  projection that drops exactly the $V(H)$-columns loses no rank on this
  framework's rigidity-row span. Because the framework is rigid on $s_c$
  and $s_c \cap V(H) = \{r\}$, the trivial motions and the projection's
  kernel together span everything, so the projected rows stay
  independent, of size at least $D(|s_c|-1)$: KT's bottom-right block,
  the pin-a-body content (\cref{lem:rank-delete-vertex}) at the row
  level.

  Third, the witness placement is KT's degenerate seed (eq.~(6.7)): the
  collapsed normal field $n \circ f$ on $G \setminus E(H)$. A per-edge
  identity --- both framings read the same support extensor, the
  projection reconciling their differing endpoints --- carries the
  projected independence of the relabelled contraction's rows back to
  the projected independence of the parent's surviving-edge rows at this
  seed.

  This discharges what KT \S5.1/Claim~6.4 states under the simplifying
  assumption that the panel coordinates are algebraically independent
  over $\Q$, realizing that genericity as the contraction's inductive
  general-position realization, the transport as the endpoint-selector
  agreement, and the rank preservation as the zero-rank-loss of the
  exterior-column projection on a rigid block.
\end{proof}

\begin{lemma}[Case I realization: simple-contraction case, all deficiencies; KT Lemma~6.3]
  \label{lem:case-I-realization}
  \lean{CombinatorialRigidity.Molecular.PanelHingeFramework.case_I_realization_all_k_gen}
  \leanok
  \uses{def:rank-hypothesis, lem:case-I, lem:rigidContract-isMinimalKDof,
        lem:rank-polynomial-IH-relabel-proj, lem:genericity-device,
        lem:rows-polynomial-in-normals, lem:isInfRigidOn-of-relative-count}
  Let $G$ be a simple minimal $k$-dof-graph with $|V(G)| \ge 3$, a proper
  rigid subgraph $H$ on $s = V(H)$ with representative body $r$, such
  that the contraction $G/E(H)$ is itself simple (KT Lemma~6.3's case
  hypothesis; the non-simple-contraction case is
  \cref{lem:case-I-realization-h65}). If the induction hypothesis
  supplies generic realizations of both $H$ (at deficiency $0$, since $H$
  is rigid) and the contraction $G/E(H)$ (at deficiency $k$, minimal
  $k$-dof by \cref{lem:rigidContract-isMinimalKDof}), then $G$ carries a
  generic realization at rank $D(|V(G)|-1) - k$ (\cref{def:rank-hypothesis}),
  discharging the Case~I contraction premise of \cref{thm:theorem-55} in
  this case.
\end{lemma}
\begin{proof}
  \leanok
  \uses{lem:rows-polynomial-in-normals, lem:rank-polynomial-IH-relabel-proj,
        lem:genericity-device, lem:case-I, lem:rigidContract-isMinimalKDof,
        lem:isInfRigidOn-of-relative-count}
  Work on a single framework $F$ at a general-position seed $q_0$, KT's
  block-triangular rank addition (eq.~(6.3)). The rigid-block rows
  ($E(H)$) are independent and number $D(|V(H)|-1)$, from the $H$-leg's
  inductive generic realization at deficiency $0$, fed through the row
  coordinatization of \cref{lem:rows-polynomial-in-normals}. The
  surviving-edge rows ($E(G) \setminus E(H)$, on $s_c = (V(G) \setminus
  V(H)) \cup \{r\}$) number $D(|s_c|-1) - k$ and, after the
  exterior-column projection killing the $V(H)$ columns, are independent
  (\cref{lem:rank-polynomial-IH-relabel-proj}, the deficiency-aware
  surviving-block rank polynomial applied to the contraction's inductive
  realization). The rigid-block rows vanish under this projection (KT's
  top-right zero block), so the two row families are disjoint and their
  union gives $D(|V(H)|-1) + D(|s_c|-1) - k = D(|V(G)|-1) - k$
  independent rows (using $|s_H \cap s_c| = 1$). The genericity device
  (\cref{lem:genericity-device}) then reads rigidity on $V(G)$ off this
  row count via the relative-count adapter
  (\cref{lem:isInfRigidOn-of-relative-count}), at $q_0$ itself; since
  $q_0$ is general position, the realization is generic.
\end{proof}

\begin{lemma}[Case I realization: non-simple $G$; KT Lemma~6.2]
  \label{lem:case-I-realization-nonsimple}
  \lean{CombinatorialRigidity.Molecular.case_I_realization_nonsimple}
  \leanok
  \uses{def:genuine-hinge-realization, lem:case-I, lem:two-vertex-zero-dof,
        lem:rigidContract-isMinimalKDof}
  Let $G$ be a minimal $k$-dof-graph with $|V(G)| \ge 3$ that is
  \emph{not} simple: it carries a parallel pair, edges $e \ne f$ both
  linking vertices $a \ne b$. Given the induction hypothesis for every
  minimal $k'$-dof-graph with fewer vertices, $G$ has a nondegenerate-hinge
  panel realization at rank $D(|V(G)|-1) - k$
  (\cref{def:genuine-hinge-realization}).
\end{lemma}
\begin{proof}
  \leanok
  \uses{lem:extensor-pair-in-panel, lem:theorem-55-base,
        lem:trivial-motions-rank-bound, lem:rigidityRows-splice-rank-add,
        lem:extProj-preserves-rank-of-inter}
  The parallel pair $\{e,f\}$ spans a proper rigid subgraph
  $H' = G[\{a,b\}] \restriction \{e,f\}$, $0$-dof by
  \cref{lem:two-vertex-zero-dof} and proper since $|V(G)| \ge 3$.
  Contract to $G' = G/H'$ and apply the induction hypothesis to obtain a
  panel-hinge framework $F_c$ on $G'$; install coincident panels at $a$
  and $b$ by pulling back $F_c$'s normal at the contracted body, and
  choose two linearly independent extensors $C_e, C_f$ in that shared
  panel (\cref{lem:extensor-pair-in-panel}). The two-body base case
  (\cref{lem:theorem-55-base}) makes the $H'$-block rigid at rank $D$;
  projecting away the $H'$-columns, the $H'$-rows vanish while $F_c$'s
  rows inject (\cref{lem:extProj-preserves-rank-of-inter}), so the two
  blocks are jointly independent and $D + \operatorname{rank}(F_c) \le
  \operatorname{rank}(F)$ (\cref{lem:rigidityRows-splice-rank-add}). The
  upper bound $\operatorname{rank}(F) \le D(|V(G)|-1)-k$ follows from the
  codimension bound (\cref{lem:trivial-motions-rank-bound}); the
  arithmetic $D + (D(|V(G)|-2)-k) = D(|V(G)|-1)-k$ closes rank equality.
\end{proof}

\begin{lemma}[Case I realization: removing a degree-two vertex; KT Lemma~6.5, Claim~6.6]
  \label{lem:case-I-realization-h65}
  \lean{CombinatorialRigidity.Molecular.PanelHingeFramework.case_I_realization_h65_gen}
  \leanok
  \uses{def:rank-hypothesis, lem:genericity-device}
  Let $G$ be a simple minimal $0$-dof-graph with $|V(G)| \ge 3$ in which
  every proper rigid subgraph has a non-simple contraction. Given the
  induction hypothesis for every smaller minimal $0$-dof-graph, $G$
  carries a generic full-rank realization (\cref{def:rank-hypothesis}).
\end{lemma}
\begin{proof}
  \leanok
  \uses{lem:genericity-device}
  By KT Claim~6.6, a vertex-inclusionwise maximal proper rigid subgraph
  $G'$ of $G$ has a non-simple contraction only because some vertex
  $v \notin V(G')$ has degree two, with two neighbours $a, b \in V(G')$
  and $G = G' + v$; then $G - v$ is a smaller simple minimal $0$-dof-graph.
  The induction hypothesis gives a generic full-rank realization of
  $G - v$ with algebraically independent seed $q$. Re-attach $v$ at the
  same seed: keep every normal and endpoint from $G-v$'s realization,
  give $v$ its own normal $q(v)$, and direct the two new edges $va, vb$
  from $v$. Off $\{va, vb\}$ the framework agrees with the inductive
  realization, so its $D(|V(G-v)|-1)$ independent rigidity rows transport
  unchanged, vanishing on $v$'s screw column. The two new rows at $va,
  vb$ span the full $D$-dimensional block through $v$'s column: their
  supporting extensors $C(q(v), q(a))$ and $C(q(v), q(b))$ are linearly
  independent because the triple $(q(v), q(a), q(b))$ is --- $q$ being
  algebraically independent and $G$ having only the three vertices $v,
  a, b$ forces this by a direct computation, with no perturbation
  argument needed. The two row blocks are disjoint (one vanishes on,
  the other is pinned through, $v$'s column) and together give $D +
  D(|V(G-v)|-1) = D(|V(G)|-1)$ independent rows, so $G$ is
  infinitesimally rigid on $V(G)$ at this seed; the genericity device
  (\cref{lem:genericity-device}) upgrades this fixed-seed rigidity to a
  generic realization.
\end{proof}

\begin{lemma}[Case I dispatch: KT Lemmas 6.2, 6.3, 6.5]
  \label{lem:case-I-dispatch}
  \lean{CombinatorialRigidity.Molecular.case_I_dispatch_gen}
  \leanok
  \uses{def:genuine-hinge-realization, def:rank-hypothesis,
        lem:case-I-realization-nonsimple, lem:case-I-realization,
        lem:case-I-realization-h65}
  Let $G$ be a minimal $0$-dof-graph with $|V(G)| \ge 3$ carrying a
  proper rigid subgraph. Then $G$ has a nondegenerate-hinge panel
  realization at the target rank (\cref{def:genuine-hinge-realization}),
  and if $G$ is simple, a generic one (\cref{def:rank-hypothesis}).
\end{lemma}
\begin{proof}
  \leanok
  \uses{lem:case-I-realization-nonsimple, lem:case-I-realization,
        lem:case-I-realization-h65, lem:generic-yields-genuine-hinge}
  If $G$ is not simple, \cref{lem:case-I-realization-nonsimple} (KT
  Lemma~6.2) applies directly. If $G$ is simple, case on whether some
  proper rigid subgraph $H$ with representative body $r$ has a simple
  contraction $G/E(H)$: if so, \cref{lem:case-I-realization} (KT
  Lemma~6.3) gives the generic realization; otherwise every contraction is
  non-simple and \cref{lem:case-I-realization-h65} (KT Lemma~6.5) gives
  it instead. Either way \cref{lem:generic-yields-genuine-hinge} then
  gives the nondegenerate-hinge form.
\end{proof}
